% !TeX root = ../sections/03_solutions.tex
\subsection{対策-周期関数}
\begin{exercise}[周期関数の分類]
	周期 \(2\pi\) をもち，区間 \([-\pi,\pi)\) 上で
	\[
	\frac{1}{x},\quad \sqrt{|x|},\quad x,\quad x^2,\quad \sin 3x
	\]
	で与えられる五つの周期関数について，次の条件に当てはまるものを全て挙げよ．
	
	\begin{enumerate}
		\item 連続な関数
		\item 区分的に連続な関数
		\item 滑らかな関数
		\item 区分的に滑らかな関数
	\end{enumerate}
\end{exercise}

\begin{background}
	この問題では，まず「区間 \([-\pi,\pi)\) 上の式」だけを見るのではなく，
	それを周期 \(2\pi\) で周期拡張した関数として考えることが重要である。
	
	周期拡張した関数を調べるときは，次の二種類の点に注意する。
	
	\begin{enumerate}
		\item 区間の内部の特異点
		\item 周期的につなぎ合わせる端点
	\end{enumerate}
	
	具体的には，区間 \([-\pi,\pi)\) 上で与えられた関数を
	周期 \(2\pi\) で延長すると，\(-\pi\) と \(\pi\) が同じ点としてつながる。
	したがって，連続性や滑らかさを判定するときは，
	\[
	x=-\pi,\quad x=\pi
	\]
	のつながり方も確認しなければならない。
	
	\medskip
	
	連続な関数とは，周期拡張したあとで全ての点において連続な関数である。
	特に，端点で値がつながっているかを確認する必要がある。
	
	区分的に連続な関数とは，有限個の不連続点を除いて連続であり，
	各不連続点で左右極限が有限に存在するような関数である。
	したがって，ジャンプ型の不連続は許されるが，
	\[
	\frac{1}{x}
	\]
	のように無限大に発散する特異点は許されない。
	
	滑らかな関数とは，ここでは周期拡張した関数が十分何回でも微分可能な関数と考える。
	したがって，関数の値だけでなく，導関数も端点でつながっている必要がある。
	
	区分的に滑らかな関数とは，有限個の点を除いて滑らかであり，
	各区間で導関数がよく振る舞う関数である。
	典型的には，有限個のジャンプや折れ曲がりは許されるが，
	導関数が無限大に発散するような点には注意が必要である。
\end{background}
\begin{strategy}
	分類問題では，次の順番で判定するとよい。
	
	\begin{enumerate}
		\item まず，関数がそもそも有限の値を持つかを確認する。
		\item 次に，区間内部で連続かどうかを見る。
		\item その後，周期拡張による端点のつながりを見る。
		\item 最後に，導関数まで見て滑らかさを判定する。
	\end{enumerate}
	
	各関数については，次のように見る。
	
	\[
	\frac{1}{x}
	\]
	は \(x=0\) で定義できず，さらに \(x\to 0\) で無限大に発散する。
	したがって，区分的に連続な関数にもならない。
	
	\[
	\sqrt{|x|}
	\]
	は関数そのものは連続である。
	しかし，\(x=0\) で尖っており，導関数
	\[
	\frac{d}{dx}\sqrt{|x|}
	\]
	が \(x=0\) の近くで無限大に発散する。
	したがって，連続ではあるが，滑らかではない。
	
	\[
	x
	\]
	は区間内部では滑らかである。
	しかし，周期拡張すると
	\[
	\lim_{x\to \pi-0} x=\pi,\qquad f(-\pi)=-\pi
	\]
	となるため，端点でジャンプが生じる。
	したがって，連続ではないが，区分的に連続であり，区分的に滑らかである。
	
	\[
	x^2
	\]
	は端点で
	\[
	(-\pi)^2=\pi^2
	\]
	となるため，関数の値はつながる。
	したがって連続である。
	しかし導関数は
	\[
	\frac{d}{dx}x^2=2x
	\]
	であり，端点で
	\[
	\lim_{x\to \pi-0}2x=2\pi,\qquad
	\lim_{x\to -\pi+0}2x=-2\pi
	\]
	となって一致しない。
	したがって滑らかではないが，区分的に滑らかである。
	
	\[
	\sin 3x
	\]
	は三角関数であり，周期拡張しても値も導関数も自然につながる。
	したがって滑らかな関数である。
	
	\medskip
	
	以上より，典型的には次のように分類できる。
	
	\[
	\begin{array}{c|l}
		\text{分類} & \text{該当する関数} \\ \hline
		\text{連続な関数}
		& \sqrt{|x|},\ x^2,\ \sin 3x \\
		\text{区分的に連続な関数}
		& \sqrt{|x|},\ x,\ x^2,\ \sin 3x \\
		\text{滑らかな関数}
		& \sin 3x \\
		\text{区分的に滑らかな関数}
		& x,\ x^2,\ \sin 3x
	\end{array}
	\]
	
	ただし，授業での「区分的に滑らか」の定義が，
	単に有限個の点を除いて各区間で滑らかであることだけを要求する場合には，
	\(\sqrt{|x|}\) の扱いが変わる可能性がある。
	この問題では，導関数の発散を考慮して，
	\(\sqrt{|x|}\) は区分的に滑らかには含めない方針で整理しておく。
\end{strategy}